\newcommand{\auditoriaid}{A00 -- Resumen consolidado}
% Preambulo comun de las auditorias del informe E1 (revision_e1/)
\documentclass[11pt,a4paper]{article}
\usepackage[utf8]{inputenc}
\usepackage[T1]{fontenc}
\usepackage[spanish,es-noshorthands]{babel}
\usepackage{lmodern}
\usepackage{microtype}
% Sin patrones de guionado en espanol en esta instalacion (ver A10): se
% desactiva el guionado automatico para evitar cortes con reglas del ingles.
\hyphenpenalty=10000
\exhyphenpenalty=10000
\usepackage[a4paper,margin=2.2cm]{geometry}
\usepackage{booktabs,array,longtable}
\usepackage{graphicx,caption,subcaption}
\usepackage{xcolor}
\usepackage{enumitem}
\usepackage{fancyhdr}
\usepackage{amsmath}
\usepackage{tikz}
\usetikzlibrary{arrows.meta,positioning}
\usepackage{natbib}
\usepackage{xurl}
\usepackage[hidelinks]{hyperref}
\urlstyle{tt}
\newcommand{\archivo}[1]{\url{#1}}
\newcolumntype{L}[1]{>{\raggedright\arraybackslash}p{#1}}
\definecolor{alta}{RGB}{180,30,30}
\definecolor{media}{RGB}{200,120,0}
\definecolor{baja}{RGB}{60,110,160}
\definecolor{cajafondo}{RGB}{245,247,250}
\newcommand{\sevA}{\textcolor{alta}{\textbf{Alta}}}
\newcommand{\sevM}{\textcolor{media}{\textbf{Media}}}
\newcommand{\sevB}{\textcolor{baja}{\textbf{Baja}}}
\setlength{\parskip}{0.35em}
\setlength{\emergencystretch}{3em}
\setlength{\parindent}{0pt}
\pagestyle{fancy}
\fancyhf{}
\fancyhead[L]{\small Auditoría del Informe del Entregable~1 -- \auditoriaid}
\fancyhead[R]{\small \thepage}
\renewcommand{\headrulewidth}{0.4pt}
% Entorno de hallazgos: ID | Severidad | Ubicacion | Hallazgo y evidencia | Correccion
\newenvironment{hallazgos}{%
\small
\begin{longtable}{@{}L{1.15cm}L{1.25cm}L{1.9cm}L{6.0cm}L{\dimexpr\textwidth-10.3cm-10.6\tabcolsep\relax}@{}}
\toprule
\textbf{ID} & \textbf{Sev.} & \textbf{Ubicación} & \textbf{Hallazgo y evidencia} & \textbf{Corrección aplicada} \\
\midrule\endfirsthead
\toprule
\textbf{ID} & \textbf{Sev.} & \textbf{Ubicación} & \textbf{Hallazgo y evidencia} & \textbf{Corrección aplicada} \\
\midrule\endhead
\bottomrule\endfoot
}{\end{longtable}}
% Macros compartidas con el informe revisado (las secciones propuestas las usan)
% Macros compartidas entre el informe revisado y las auditorias.
% Cifras verificadas contra los archivos del repositorio (ver A01/A04/A07):
\newcommand{\nUniverso}{102}   % source_id CMIP6 con tos/Omon en ESGF (00b, model_availability_priority.csv)
\newcommand{\nSeed}{41}        % publican los cuatro experimentos (config/models_seed_cmip6.csv)
\newcommand{\nSel}{40}         % completos + QC PASS (models_inventory_final.csv)
\newcommand{\nComb}{160}       % combinaciones modelo-experimento PASS (qc_report.csv)
\newcommand{\nNoEnc}{44}       % config/models_no_encontrado_44.csv
\newcommand{\nParcial}{18}     % 102 - 40 - 44
% Fecha de la portada: fija (no \today), editable en un solo lugar.
\newcommand{\fechaentrega}{28 de septiembre de 2026}
% Estado de codigo que documenta el E1 (antes de los cambios del E2)
\newcommand{\commitEone}{\texttt{8ee6184}}

% En las auditorias el texto propuesto se muestra fuera del informe: las
% referencias cruzadas se imprimen con el nombre de la seccion destino.
\makeatletter
\def\etq#1#2{\expandafter\def\csname etq@#1\endcsname{#2}}
\etq{sec:resumen}{Resumen}
\etq{sec:alcance}{Alcance}
\etq{sec:marco}{Revisión científica}
\etq{sec:antecedentes}{Antecedentes}
\etq{sec:referencias}{Literatura sobre TOE}
\etq{tab:referencias}{bibliografía TOE}
\etq{sec:similitudes}{Niño~3.4 frente a Niño~1+2}
\etq{sec:toe_metodos}{Elección de los métodos de TOE}
\etq{tab:metodos_toe}{métodos TOE}
\etq{eq:szabo_signal}{ecuación señal M1}
\etq{eq:szabo_V}{ecuación V(t)}
\etq{eq:szabo_T}{ecuación T(t)}
\etq{eq:szabo_toe}{ecuación TOE M1}
\etq{eq:tod}{ecuación ToD}
\etq{sec:datos}{Datos}
\etq{fig:cmip6_contribution}{contribución institucional}
\etq{sec:tabla_modelos}{Modelos seleccionados}
\etq{tab:modelos}{modelos}
\etq{sec:desviacion}{Selección de modelos}
\etq{sec:dominios_periodos}{Dominios y periodos}
\etq{fig:region_nino}{regiones Niño}
\etq{sec:arquitectura}{Arquitectura del pipeline}
\etq{fig:flujograma}{Flujograma}
\etq{sec:scripts}{Descripción de los pasos}
\etq{sec:paso00b}{Paso 00b}
\etq{sec:paso01}{Paso 01}
\etq{sec:paso02}{Paso 02}
\etq{sec:paso04}{Paso 04}
\etq{sec:paso05}{Paso 05}
\etq{sec:paso06}{Paso 06}
\etq{sec:paso07}{Paso 07}
\etq{sec:reporte}{Reporte de estado}
\etq{sec:resultados}{Resultados}
\etq{fig:qc}{matriz de estado}
\etq{tab:cuentas}{conteo de modelos}
\etq{fig:mapas2d}{mapas 2D}
\etq{fig:series}{series}
\etq{fig:boxplot}{boxplots}
\etq{sec:roadmap}{Próximos pasos}
\etq{tab:roadmap}{próximos pasos}
\etq{sec:conclusiones}{Conclusiones}
\etq{sec:recomendaciones}{Recomendaciones}
\makeatother

\makeatletter
\AtBeginDocument{\renewcommand{\ref}[1]{\@ifundefined{etq@#1}{\textit{[#1]}}{\textit{\csname etq@#1\endcsname}}}}
\makeatother
% Texto propuesto: secciones degradadas un nivel y en caja
\newenvironment{propuesto}{%
  \par\medskip\noindent\textcolor{baja}{\rule{\textwidth}{0.8pt}}\par
  \noindent{\small\textbf{Inicio del texto propuesto para el informe revisado}}\par\medskip
  \begingroup
  \let\section\subsection\let\subsection\subsubsection\let\subsubsection\paragraph
}{%
  \endgroup
  \par\noindent{\small\textbf{Fin del texto propuesto}}\par
  \noindent\textcolor{baja}{\rule{\textwidth}{0.8pt}}\par\medskip}
% Recuadro de actualizacion posterior (decisiones del autor)
\newcommand{\actualizacion}[1]{\par\medskip\noindent\fcolorbox{media}{media!8}{\parbox{\dimexpr\linewidth-2\fboxsep-2\fboxrule\relax}{\small\textbf{Actualización del 29-09-2026 (decisiones del autor).} #1}}\par\medskip}

\begin{document}

\begin{center}
{\Large\bfseries Revisión del Informe del Entregable~1\\[0.2em] Resumen consolidado de auditorías}\\[0.5em]
{\normalsize Informe original: \texttt{informe/informe\_e1\_smnv3.tex} (commit \texttt{330fd2e})}\\[0.2em]
{\small Revisión: \fechaentrega}
\end{center}

\section{Qué se hizo}

Se auditó por separado cada tema del informe del Entregable~1 (A01--A10). Cada afirmación se contrastó con cuatro fuentes: el texto del TdR; el texto completo de los artículos citados (colección local de PDF del proyecto y Crossref); el código en el estado del E1 (commit \texttt{8ee6184}); y los datos y figuras producidos. Con el texto corregido de cada auditoría se armó un informe nuevo con la misma estructura, extensión e información equivalentes, y el uso de \texttt{run.sh} pasó a un documento de anexos separado. No se modificó ningún archivo versionado del repositorio y no se hizo \emph{commit} ni \emph{push}: todo está en \archivo{revision_e1/}, carpeta excluida de git mediante \archivo{.git/info/exclude}.

\actualizacion{Por decisión del autor: (1) la Sección~6 del informe se redujo a un párrafo con las acciones para obtener los datos y el flujograma, y la descripción paso a paso pasó al anexo; (2) el informe describe la ventana de descarga vigente, 0°--360° y 30°S--30°N (ampliada para RONI y para el índice $W$ del ToD), con los mapas 2D y el mapa de regiones regenerados para toda esa ventana; (3) el techo del control de calidad es 45\,°C, justificado por el mar Rojo y el golfo Pérsico; (4) el informe no cita commits; (5) el repositorio es de libre acceso, y así se indica en el anexo. Se actualizaron el informe, el anexo y las auditorías A01, A03--A08 y A10.}

\section{Documentos entregados}

\begin{center}\small
\begin{tabular}{@{}lL{9.2cm}r@{}}
\toprule
Archivo & Contenido & Págs. \\
\midrule
\csname @@input\endcsname ../trabajo/tabla_paginas.tex 
\bottomrule
\end{tabular}
\end{center}

\section{Hallazgos por auditoría}

\begin{center}\small
\begin{tabular}{@{}lrrrr@{}}
\toprule
Auditoría & Alta & Media & Baja & Total \\
\midrule
A01 Portada, Resumen, Alcance & 2 & 5 & 3 & 10 \\
A02 Antecedentes y literatura & 3 & 9 & 3 & 15 \\
A03 Niño~3.4/1+2 y métodos TOE & 5 & 10 & 3 & 18 \\
A04 Datos & 5 & 3 & 3 & 11 \\
A05 Arquitectura y \texttt{run.sh} & 0 & 5 & 2 & 7 \\
A06 Pasos 00--07 & 0 & 8 & 4 & 12 \\
A07 Resultados y figuras & 4 & 4 & 3 & 11 \\
A08 Cierre & 2 & 6 & 2 & 10 \\
A09 Bibliografía & 0 & 5 & 3 & 8 \\
A10 Compilación y cierre & 2 & 5 & 4 & 11 \\
\midrule
\textbf{Total} & \textbf{23} & \textbf{60} & \textbf{30} & \textbf{113} \\
\bottomrule
\end{tabular}
\end{center}

Severidad \emph{alta}: afirmación falsa, contradicción con los datos o el código, o defecto visible que un revisor externo detectaría. \emph{Media}: imprecisión, omisión relevante o inconsistencia interna. \emph{Baja}: estilo y forma.

\section{Correcciones de mayor impacto}

\begin{enumerate}[leftmargin=1.5em]
\item \textbf{Composición del universo de modelos} (A01, A04). ``46 del barrido más 56 de la literatura'' era un resto de versiones anteriores. Los 102 son todo el universo ESGF con \texttt{tos}/\texttt{Omon}: 49 publican los tres experimentos del TdR, 41 publican además \texttt{ssp370}, y 40 quedan seleccionados (44 + 18 excluidos).
\item \textbf{Método M1 mal atribuido} (A02, A03). La suma de varianzas $T=V+M+S$ viene de Hawkins y Sutton (2009), no de Szabó; el contraste ``aditiva frente a raíz de suma de cuadrados'' no existe (son la misma magnitud); el umbral $\kappa=2$ es de Hawkins y Sutton (2012), no de Szabó. Se agregaron las definiciones de $G$, $M$, $S$ y $W$, que faltaban.
\item \textbf{Afirmaciones sobre la literatura} (A03). La caja EEP no incluye Niño~1+2; el análisis de Larson y McMonigal es del periodo histórico, no del siglo~XXI; Poul et al. atribuyen el retraso a menos señal, no a más ruido; Bruno evaluó 53 modelos (32 de ellos están entre los seleccionados), no 54 (37).
\item \textbf{Figuras que no correspondían al E1} (A04, A07). Por los enlaces simbólicos a \texttt{figures/}, el PDF mostraba la ventana del E2, un \emph{boxplot} viejo con 31 modelos y un periodo ``1900--2015'' causado por un error de redondeo. La figura de contribución era una imagen externa con países que no están en el conjunto. Todas se regeneraron o reemplazaron.
\item \textbf{Hoja de ruta desalineada con el E2} (A08). Anunciaba la corrección Linear Scaling, que el E2 descartó con fundamento.
\item \textbf{Ventana y reproducibilidad} (A01, A04, A05, A08, A10). El informe describía una ventana distinta de la que muestran los datos actuales. Ahora describe la ventana vigente y el techo de 45\,°C, sin citar commits.
\item \textbf{Formato} (A10). Fuentes de mapa de bits (Type~3), guionado con reglas del inglés, cuatro pasadas necesarias y ancla \texttt{page.1} duplicada.
\end{enumerate}

\section{Resultados nuevos verificados que se incorporaron}

\begin{itemize}[leftmargin=1.2em]
\item Variabilidad observada (ERSSTv5, 1981--2014): la desviación estándar de las anomalías mensuales es de 1{,}08\,°C en Niño~1+2 frente a 0{,}86\,°C en Niño~3.4. Da sustento propio a la hipótesis de emergencia más tardía en Niño~1+2 (A03).
\item Sesgo de la mediana histórica (1900--2014): Niño~1+2 es más cálido que ERSSTv5 en 38 de 40 modelos (sesgo mediano de $+2{,}2$\,°C); Niño~3.4 tiene un sesgo mediano de $-0{,}24$\,°C. Es coherente con Erickson y Patricola (2023), que ahora se cita (A07).
\item Costo de exigir \texttt{ssp370}: 8 modelos con los tres experimentos del TdR quedan fuera (49 → 41). Se agregó como recomendación (A04, A08).
\item Cobertura temporal verificada archivo por archivo: los 40 \texttt{historical} terminan en 2014, sin huecos. Las únicas excepciones de periodo son IITM-ESM y CAMS-CSM1-0 (A06).
\end{itemize}

\section{Decisiones que quedan para el autor}

\begin{enumerate}[leftmargin=1.5em]
\item \textbf{Fecha de la portada}: el informe y el anexo usan ``\fechaentrega'' (macro \verb|\fechaentrega| en \archivo{revision_e1/trabajo/macros_informe.tex}).
\item \textbf{Acciones en el repositorio} (A10, sección~3): congelar las figuras como archivos reales; corregir los scripts de gráficos; actualizar comentarios y CSV desactualizados; adoptar la bibliografía revisada. Ninguna se ejecutó.
\item \textbf{Instalar} \texttt{texlive-lang-spanish} (requiere \texttt{sudo}) para reactivar el guionado en español.
\item \textbf{Recomendación sobre los 8 modelos sin \texttt{ssp370}}: se incluyó en el informe revisado como recomendación; puede retirarse si no se quiere abrir esa discusión en el E1.
\end{enumerate}

\section{Extensión del informe}

\begin{center}\small
\begin{tabular}{@{}lrr@{}}
\toprule
Sección & Original (palabras) & Revisado (palabras) \\
\midrule
Resumen y Alcance & 675 & 698 
 \\
Antecedentes y literatura & 972 & 1055 
 \\
Niño~3.4/1+2 y métodos TOE & 1344 & 1126 
 \\
Datos & 1599 & 1456 
 \\
Arquitectura (el uso pasa al anexo) & 749 & 307 
 \\
Obtención de datos (antes pasos 00--07) & 1628 & 479 
 \\
Resultados & 952 & 904 
 \\
Próximos pasos, conclusiones, recomendaciones & 629 & 654 
 \\
\midrule
\textbf{Total de secciones} & \textbf{8\,548} & \textbf{6\,675} \\
\bottomrule
\end{tabular}
\end{center}

Con el mismo método de conteo (\texttt{detex} por sección, tablas incluidas), las secciones originales suman 8\,548 palabras. Por decisión del autor, la descripción de los pasos (unas 1\,400 palabras) y el uso de \texttt{run.sh} (unas 600) pasaron al anexo, por lo que el informe es más corto; esa información se conserva completa en el anexo, que creció a 12 páginas.

\end{document}
